\section{Discussion}
\label{sec:discussion}

\subsection{From Single-Day to Multi-Day: A Unified Interpretation}

The progression from single-day validation to multi-day regime detection reveals a coherent picture of LLM structural reasoning capabilities. Single-day analysis establishes that LLMs detect dealer constraints through genuine mechanism identification: the raw chain superiority (92.3\% vs.\ 61.5\%) proves the model reconstructs dealer positioning from first principles, while the inverse p-hacking finding (33\% lower range expansion) and alpha orthogonality (Sharpe 1.8 $\rightarrow$ 0.1) confirm the model identifies structural states rather than exploitable anomalies.

These single-day findings provide necessary grounding for interpreting regime detection results. The 69.1pp discrimination between 2024 and 2020 is credible precisely because we have established that the LLM performs genuine structural analysis---not pattern matching on pre-calculated metrics. The alpha disappearance puzzle from single-day work motivates the regime extension: if dealers maintain persistent positioning, the structural constraint is detectable but not tradeable because it never transitions. Multi-day analysis confirms this interpretation: 100\% detection in 2024--2025 reflects sustained structural coherence rather than detection-profitable switching.

Together, both temporal scales support the same conclusion: LLMs identify emergent constraints in complex adaptive systems---coordination patterns arising from decentralized dealer inventory management~\cite{grossman1988liquidity}---without capturing the dynamic adaptations that generate trading alpha. This positions temporal obfuscation testing as a domain-agnostic methodology for validating structural reasoning wherever training data contamination threatens evaluation validity.

\subsection{Regime Exceptionality vs.\ Universal Positioning}

The selectivity finding is critical. Unlike prior 5-day testing that achieved universal 98--100\% detection, 30-day windows with strict criteria produce selective 12--81\% detection. Four empirical results support regime exceptionality over universal pattern matching:

First, 2020 fragmented markets (12.1\% detection) confirm that dealer hedging existed but lacked the persistence and magnitude characterizing 2024 regimes. Second, the gradual adoption pattern (3.7\% $\rightarrow$ 100\%) tracks a known structural variable (0DTE volume share), ruling out detection artifacts. Third, negative controls (0\% FP on transitional/low-magnitude) confirm the framework enforces criteria, not merely detecting ``negative GEX.'' Fourth, the regime sign flip (positive dominant in 2020, negative dominant in 2024) demonstrates sensitivity to directional structure.

\subsection{Alternative Explanations and Robustness}

Price normalization analysis addresses the inflation confound: applying SPY's 2.3$\times$ price growth to 2020 magnitudes accounts for only 31\% of the detection gap, with 69\% attributable to structural evolution in persistence (96\% vs.\ 69\%) and stability (99\% vs.\ 54\%).

The shuffle test (Phase 2a) warrants interpretation. The 61.1\% false positive rate on 2024 shuffled data appears concerning but is itself an important structural finding: 2024 regimes exhibit such extreme same-sign dominance (96\%) that randomizing day order rarely breaks the persistence threshold. This validates that 2024 detection reflects aggregate statistical properties robust to reordering, while 2020 detection (12.1\% on shuffled) depends on temporal structure---consistent with fragmented markets where day ordering matters.

Formula Agreement Testing confirms detection is robust to mathematical formulation, with 100\% agreement (61/61 windows) between absolute dollar-scaled GEX and normalized GEX ratios, demonstrating the model identifies regimes based on structural positioning rather than nominal magnitude cues.

\subsection{Practitioner Implications}

For practitioners deploying LLM-based market analysis, our findings suggest specific considerations:

\textbf{Role separation}: LLMs excel at mechanism identification (90.9\% accuracy) but underperform at probability calibration (AUC 0.488 vs.\ statistical model's 0.681). Optimal deployment uses LLMs for structural detection with statistical models providing probability estimates.

\textbf{Raw data advantage}: The 30.8pp improvement from raw chain data suggests providing LLMs with granular strike-level data rather than summary statistics. Net GEX may actually \textit{obscure} structural signals visible in the full options surface.

\textbf{Detection $\neq$ alpha}: The stable detection/declining profitability pattern demonstrates that structural pattern detection does not guarantee trading edge. Detected patterns serve as risk management inputs---structural constraints that suppress volatility are precisely the regime information risk models require.

\textbf{Prompt design}: The 28.5pp biased-unbiased gap reveals significant framing sensitivity. Practitioners should employ minimal prompting for core reasoning, avoiding keyword enumeration that triggers template responses.

\subsection{Limitations}

\textbf{Single asset}: Testing exclusively on SPY limits generalizability, though SPY concentrates dealer hedging in the world's most liquid market. Cross-asset validation, particularly on equities with positive gamma regimes, remains essential.

\textbf{End-of-day granularity}: Daily snapshots miss intraday dynamics. However, EOD data achieves AUC = 0.681 for T+1 outcomes, validating sufficient predictive information for overnight positioning decisions.

\textbf{Single model family}: Validation using o3-mini/o4-mini cannot distinguish architecture-specific capabilities from general LLM reasoning. Cross-model validation would strengthen generalizability.

\textbf{Raw chain sample size}: The $n = 13$ pilot provides strong initial evidence (effect exceeds minimum detectable by $2\times$) but requires larger-scale replication.

\textbf{Shuffle test limitation}: The 61\% false positive rate on 2024 shuffled data reflects extreme regime dominance rather than framework failure, but indicates the persistence criterion alone is insufficient for highly polarized markets. This is documented as a known boundary condition.

\textbf{Data gaps}: Alpha Vantage provides 4 AM -- 8 PM ET data only, precluding overnight transmission analysis to global sessions. Intraday 0DTE dynamics are captured only through end-of-day open interest residuals.
